\documentclass[11pt,letterpaper]{article}
\usepackage[margin=1in]{geometry}
\usepackage{indentfirst,amsmath,amssymb,enumitem}

\usepackage[utf8]{inputenc}
\usepackage[T2A]{fontenc}
\usepackage{textcomp}

\usepackage{xspace,xcolor}
\definecolor{green1}{HTML}{198019} % green from 1505.03893

\usepackage[
  hyperfootnotes=false,
  linktoc=all,
  colorlinks=true,
  linkcolor=green1,
  citecolor=green1,
  urlcolor=blue,
  pdfauthor={Ivan Pogrebnyak},
  pdftitle={Units in C++},
  pdfsubject={},
  pdfkeywords={units, C++}
]{hyperref}

\title{Units in C++}
\author{Ivan Pogrebnyak}
\date{\today}

\begin{document}

\maketitle

\section{Introduction} % ===========================================================================

Any physics or engineering simulation software has to account for units,
whether it is done explicitly or implicitly.
In this document, we discuss a new units library for C++ that leverages the type system to
represent units.
The library provides features to define dimensionful quantities,
perform arithmetic operations on them, and validate compatibility of quantities at compile time.

\section{Theory} % =================================================================================

A physical \textbf{quantity}, such as volume or pressure, can be divided into two components:
\textbf{units}, specifying the type of quantity, and the quantity's \textbf{value}, specifying how
large the quantity is in the given units.

Geometrically, units can be viewed as \textbf{dimensions} of quantities. Much like a point
in 3-dimensional space can be represented by a vector $(x,y,z)$, quantities in mechanics have
dimensions of mass, length, and time. The units' dimensions are discrete-valued. For example,
represented as (mass, length, and time) vectors, speed units have dimensions $(0,1,-1)$, and those
for energy are $(1,2,-2)$.

Algebraically, units are multiplicative factors, orthogonal to the quantity's value.
For example, a length quantity can be expressed as 5 m, which can be thought of as multiplication
between a numeric value 5 and units m. Two lengths multiplied yields an area,
$5~\text{m} \times 3~\text{m} = 15~\text{m}^2$.
Note that both values and units were multiplied. Addition is typically only useful to define between
quantities with the same dimensions, while multiplication and division are meaningful in all cases.
For example, $\text{force} / \text{mass} = \text{acceleration}$.

Apples and oranges don't really behave algebraically like units. One is not typically interested in
oranges$^2$, while they may want to buy $5~\text{apples} + 3~\text{oranges}$.

Typically, multiple units exist for the same type of quantity, i.e. vector of dimensions. For
example, length can be measured in m, mm, km, mi, in, ft, etc. Mathematically, conversions between
these redundant units amount to multiplication by dimensionless factors. However, floating-point
arithmetic on a computer is not associative. When implementing a units system, it is important to
think about how conversions are handled and at what point. And, as always, one must find the optimal
trade-off between numeric precision, user friendliness, difficulty of implementation and
maintenance, and overhead at run time and compile time.

\section{The units library} % =================================================================

\subsection{Example}

To illustrate what problems the library addresses and how it solves them,
let us start with a simple example.
Consider a physics simulation involving projectile motion.
Projectiles may be represented as classes with member variable containing relevant quantities,
such as projectile's mass, position, velocity, etc.
Without explicitly accounting for units, the implementation may look like this.

\section{Implementation choices} % =================================================================

\subsection{Run-time or compile-time representation of units}

\subsection{Immediate or delayed conversion}

\subsection{Static or dynamic number of dimensions}

\subsection{Distinguishing units with equal dimensions but distinct meaning}

Dimensionless quantities: degrees, radians. Degrees as a factor: std::sin(15 * "deg"\_u).

\section{Future features} % ========================================================================

Features that can but haven't been implemented.

\subsection{Rational powers}

\subsection{User-defined units and relative definitions}

\InputIfFileExists{extra}{}{}

\end{document}
